% =============================================================================
% 03-kata-pengantar.tex — Kata Pengantar
% =============================================================================
\clearpage
\phantomsection
\addcontentsline{toc}{chapter}{KATA PENGANTAR}

\begin{center}
  {\fontsize{14}{16}\selectfont\bfseries KATA PENGANTAR}
\end{center}

\vspace{1cm}

Puji dan syukur penulis panjatkan kepada Tuhan Yang Maha Esa atas berkat dan rahmat-Nya sehingga penulis dapat menyelesaikan laporan Kerja Praktik yang berjudul ``\judulLaporan'' dengan baik dan tepat waktu.

Laporan ini disusun sebagai salah satu syarat untuk menyelesaikan mata kuliah Kerja Praktik pada Program Studi Informatika, Fakultas Teknologi Cerdas, Universitas Kristen Krida Wacana. Kerja Praktik ini dilaksanakan menggunakan skema Proyek Dosen di bawah bimbingan langsung Ibu Astin selaku Ketua Program Studi Psikologi dan Kepala Riset \textit{Center for Health Psychology} (CHP) UKRIDA.

Penulis menyadari bahwa laporan ini tidak dapat diselesaikan tanpa bantuan dan dukungan dari berbagai pihak. Oleh karena itu, penulis ingin menyampaikan ucapan terima kasih yang sebesar-besarnya kepada:

\begin{enumerate}[leftmargin=2cm]
  \item Ibu Astin selaku pemilik proyek dan Kepala Riset CHP UKRIDA, yang telah memberikan kesempatan, arahan, dan umpan balik yang berharga selama pelaksanaan Kerja Praktik.
  \item Dosen Pembimbing Kerja Praktik yang telah membimbing penulis dalam penyusunan laporan ini.
  \item Seluruh dosen dan staf Program Studi Informatika, FTC, UKRIDA yang telah memberikan bekal ilmu dan dukungan selama perkuliahan.
  \item Rekan satu tim yang telah berkolaborasi dalam pengembangan Platform Asesmen Digital CHP.
  \item Keluarga tercinta yang senantiasa memberikan doa, motivasi, dan dukungan moral selama proses pelaksanaan Kerja Praktik.
\end{enumerate}

Penulis menyadari bahwa laporan ini masih jauh dari kesempurnaan. Oleh karena itu, penulis mengharapkan kritik dan saran yang membangun demi perbaikan di masa mendatang. Semoga laporan ini dapat memberikan manfaat bagi pembaca dan menjadi kontribusi yang berarti bagi pengembangan platform asesmen psikologi digital di Indonesia.

\vspace{2cm}

\noindent
\begin{flushright}
  Jakarta, 30 Juni \tahunSidang\\[1.5cm]
  \namaMahasiswa
\end{flushright}
